\section{Analytical analysis}\label{sec:analysis}
This section works through the pitching moment --- thrust on or off, thrust axis aligned or
misaligned --- and the static and dynamic stability results that follow. An acceleration component
such as $\vec{a}_r\cdot\hat{x}_r$ is the acceleration resolved along the body axes at the current
instant, not a derivative taken inside the rotating frame $\mathcal{F}_r$.

Two inputs are zero by default and reduce everything below to its simplest form: the thrust angle
$\theta_t$ and the gravity-direction rotation $\varphi=x/R_e$, which stays zero if the local
vertical is frozen at the launch point.

\subsection{Forces}\label{sec:analysis_forces}
Aerodynamic force and thrust are already body-axis quantities, so only weight has to be resolved.
Weight acts along the local vertical $\hat{z}_g(\varphi)=-\sin\varphi\,\hat{x}_0+\cos\varphi\,
\hat{z}_0$, so with \autoref{eq:body_axes_in_ground}
\begin{equation}
    \vec{W}\cdot\hat{x}_r = -mg\sin(\theta_r+\varphi), \qquad
    \vec{W}\cdot\hat{z}_r = +mg\cos(\theta_r+\varphi).
\end{equation}
The downrange rotation of gravity thus enters only as a shift of the pitch angle,
$\theta_r\to\theta_r+\varphi$: the pitch above the \emph{local} horizon. Adding thrust and the
aerodynamic force,
\begin{align}
    m\,\vec{a}_r\cdot\hat{x}_r &= -mg\sin(\theta_r+\varphi) + T\cos\theta_t - qSC_x \label{eq:analytic_force_x}\\
    m\,\vec{a}_r\cdot\hat{z}_r &= mg\cos(\theta_r+\varphi) - T\sin\theta_t - qSC_z \label{eq:analytic_force_z}
\end{align}
Sign checks, with $\varphi=\theta_t=0$. Nose straight up ($\theta_r=90^\circ$): weight is purely
axial, $m\,\vec{a}_r\cdot\hat{x}_r=-mg+T-qSC_x$. Nose horizontal ($\theta_r=0$): weight is purely
normal, $m\,\vec{a}_r\cdot\hat{z}_r=mg-qSC_z$. For $\varphi$, flying level at downrange $x>0$ gives
$m\,\vec{a}_r\cdot\hat{x}_r=-mg\sin\varphi<0$: the Earth's centre lies below and behind the
rocket, so the pull is tilted against the flight direction. At $x=10\,\mathrm{km}$ this is
$g\varphi\approx1.5\times10^{-2}\,\mathrm{m/s^2}$.

\subsection{Moment}\label{sec:analysis_moment}
Substituting the arms of \autoref{eq:arms} into \autoref{eq:eom_rot}, with $\hat{x}_r\times
\hat{z}_r=-\hat{y}_r$ and the thrust vector of \autoref{eq:thrust_frame}:
\begin{equation} \label{eq:moment_full}
    I\varepsilon_r = (r_g - r_t)T\cos\theta_t + (l_g - l_t)T\sin\theta_t
                     - r_g\,qSC_x + (l_g - l_p)\,qSC_z - qSc\,C_m
\end{equation}
Thrust acts at $t$, displaced from the CoG by $\vec{gt}=(l_g-l_t)\hat{x}_r+(r_g-r_t)\hat{z}_r$:
the radial offset gives an arm to the axial thrust component $T\cos\theta_t$, the axial offset an
arm to the normal component $T\sin\theta_t$. The aerodynamic force acts at $p$: the axial arm
$l_g-l_p$ carries $qSC_z$ and the radial arm $r_g$ carries $qSC_x$. Weight does not appear, since
$\vec{gg}\times\vec{W}=0$ for every orientation of $\vec{W}$. Rotating the local vertical is
therefore a purely translational correction, like the Coriolis term.

For on-axis thrust, $\vec{T}=T\hat{x}_r$, sliding the application point along the axis changes
nothing: the force is parallel to the displacement. Only $r_t$ survives. Once $\theta_t\neq0$ the
normal component does care how far along the axis it acts, hence the second term.

\autoref{eq:moment_full} specialises four ways; \emph{aligned} means $r_g=r_t=\theta_t=0$.
\begin{center}
\begin{tabular}{l l l}
 & \textbf{aligned} ($\theta_t=0$) & \textbf{misaligned} ($\theta_t\neq0$) \\
\textbf{no thrust} ($T=0$) & \autoref{eq:moment_nothrust} & identical --- $T=0$ removes $\theta_t$ \\
\textbf{with thrust} ($T>0$) & \autoref{eq:moment_aligned} & \autoref{eq:moment_misaligned} \\
\end{tabular}
\end{center}

\subsubsection{No thrust: coast and descent}\label{sec:analysis_nothrust}
Setting $T=0$ leaves the airframe's purely aerodynamic behaviour, independent of the motor:
\begin{equation} \label{eq:moment_nothrust}
    I\varepsilon_r = -r_g\,qSC_x + (l_g - l_p)\,qSC_z - qSc\,C_m.
\end{equation}
Since $T=0$ removes both thrust terms, misalignment is irrelevant here, and the static-stability
result below, built on this case, needs no revision when it is switched on.

\subsubsection{With thrust, aligned}\label{sec:analysis_aligned}
With $r_g=r_t=0$ and $\theta_t=0$,
\begin{equation} \label{eq:moment_aligned}
    I\varepsilon_r = (l_g - l_p)\,qSC_z - qSc\,C_m.
\end{equation}
Thrust has dropped out exactly: an aligned rocket's attitude dynamics are the same during boost
as during coast, however large $T$ is.

\subsubsection{With thrust, misaligned}\label{sec:analysis_misaligned}
With the nozzle canted by a small fixed $\theta_t$ and the geometry otherwise on-axis,
\begin{equation} \label{eq:moment_misaligned}
    I\varepsilon_r = \underbrace{(l_g - l_t)T\sin\theta_t}_{M_{\theta_t}}
                     + (l_g - l_p)\,qSC_z - qSc\,C_m.
\end{equation}
Unlike every aerodynamic term, $M_{\theta_t}$ survives $q=0$ and $\alpha=0$: while the motor burns,
the airframe carries a disturbance torque proportional to the thrust. Its sign: $l_t>l_g$ for any
real rocket, so $l_g-l_t<0$. A positive $\theta_t$ pushes the tail towards $-\hat{z}_r$ and pitches
the nose \emph{down}, $M_{\theta_t}<0$; a negative $\theta_t$ pitches it up.

Since $\partial M_{\theta_t}/\partial\alpha=0$, the stability derivative of
\autoref{eq:dtheta_dalpha} is untouched: misalignment shifts the equilibrium without changing
whether it is stable. Setting $\varepsilon_r=0$ with the linearised coefficients of
\autoref{eq:aero_linearised} and $C_{z,0}=C_{m,0}=0$ gives the trim angle of attack,
\begin{equation} \label{eq:alpha_trim}
    \alpha_{trim} = -\frac{M_{\theta_t}}{I\,\partial\varepsilon_r/\partial\alpha}
        = \frac{-(l_g - l_t)T\sin\theta_t}
               {qS\bigl[(l_g - l_p)C_{z,\alpha} - c\,C_{m,\alpha}\bigr]}.
\end{equation}
A stable airframe has a negative denominator, so $\alpha_{trim}$ carries the sign of
$M_{\theta_t}$: the rocket flies at a small constant angle off the relative wind, whose normal
force curves the trajectory, rather than tumbling. $\alpha_{trim}\propto T/q$ grows without bound
just after ignition and in any high-altitude burn, and as the static margin tends to zero --- the
physical reason launch rails exist.

The radial offsets $r_g,r_t$ enter \autoref{eq:moment_full} on the same footing as $\theta_t$,
adding the constant offset $(r_g-r_t)T\cos\theta_t-r_g\,qSC_x$; the worked cases set them to zero
to isolate $\theta_t$. A commanded, time-varying $\theta_t(t)$ remains future work.

\subsection{Static stability}\label{sec:analysis_stability}
Static stability asks whether a disturbance in the angle of attack $\alpha$ produces a restoring
moment. Take the no-thrust aligned case, \autoref{eq:moment_nothrust} with $r_g=0$, and linearise
the coefficients,
\begin{equation} \label{eq:aero_linearised}
    C_z = C_{z,0} + \alpha\,C_{z,\alpha}, \qquad
    C_m = C_{m,0} + \alpha\,C_{m,\alpha},
\end{equation}
with slopes $C_{z,\alpha}=\partial C_z/\partial\alpha$, $C_{m,\alpha}=\partial C_m/\partial\alpha$ and
offsets $C_{z,0}$, $C_{m,0}$. This is an analysis device applied to whichever coefficients are in
force; the Linearized model (\autoref{sec:aero_coeffs}) is its form in $\alpha$, $M$ and $Re$.
Stability requires $\partial\varepsilon_r/\partial\alpha<0$: a disturbance must produce an angular
acceleration that reduces it. The derivative is
\begin{equation} \label{eq:dtheta_dalpha}
    \frac{\partial \varepsilon_r}{\partial \alpha} = \frac{qS}{I}\biggl[(l_g - l_p)C_{z,\alpha} - c\,C_{m,\alpha}\biggr],
\end{equation}
negative for $C_{z,\alpha}>0$ (the normal case) when
\begin{equation} \label{eq:static_margin}
    l_g - l_p < c\,\frac{C_{m,\alpha}}{C_{z,\alpha}}.
\end{equation}
Both lengths are measured from the nose, so $l_g-l_p<0$ exactly when $p$ sits aft of $g$ --- the
familiar ``centre of pressure behind centre of gravity'' rule, which is what
\autoref{eq:static_margin} reduces to when $p$ is the centre of pressure itself, $C_{m,\alpha}=0$.
In this report's sign of $\vec{A}_m$, a positive $C_{m,\alpha}$ is stabilising, the opposite of
the textbook $C_{m,\alpha}<0$.

\paragraph{Static margin, in calibers.} $p$ is a moment reference station, not necessarily the centre
of pressure (\autoref{sec:geometry}): the CP is the derived station at which the moment slope about it
vanishes,
\begin{equation}
    l_{cp} = l_p - c\,\frac{C_{m,\alpha}}{C_{z,\alpha}},
\end{equation}
and the rocketry convention divides the CP-to-CoG arm by the reference length, flipping the sign so a
stable rocket gets a positive number:
\begin{equation}
    \mathrm{SM} = \frac{l_{cp}-l_g}{c} = \frac{l_p-l_g}{c} - \frac{C_{m,\alpha}}{C_{z,\alpha}}.
\end{equation}
$\mathrm{SM}>0$ is then equivalent to \autoref{eq:static_margin}; the two coincide, $l_{cp}=l_p$, exactly
when $C_{m,\alpha}=0$ about $p$ (the degenerate case where $p$ already is the CP).

\subsection{Flight-path angle and the point-mass gravity turn}\label{sec:analysis_flightpath}
Everything above tracks the \emph{nose} direction $\theta_r$. The trajectory itself is set by the
\emph{velocity} direction, and the two differ by the angle of attack. Define the flight-path angle
$\theta_v$, the angle of $\vec{v}_r$ above the horizon (conventionally written $\gamma$; here
$\theta_v$, since $\gamma_r$ is the roll angle of \autoref{sec:roll}), by the same construction as
\autoref{eq:body_axes_in_ground}:
\begin{equation}\label{eq:vel_axes_in_ground}
    \hat{x}_v = \cos\theta_v\,\hat{x}_0 - \sin\theta_v\,\hat{z}_0, \qquad
    \hat{z}_v = \sin\theta_v\,\hat{x}_0 + \cos\theta_v\,\hat{z}_0, \qquad
    \vec{v}_r = V\hat{x}_v,
\end{equation}
so that $v_x=V\cos\theta_v$, $v_z=-V\sin\theta_v$ and $\theta_v=\operatorname{atan2}(-v_z,v_x)$:
positive climbing, since $\hat{z}_0$ points down. In the pitch plane and in still air the three
angles close exactly,
\begin{equation}\label{eq:alpha_closure}
    \alpha = \theta_r - \theta_v,
\end{equation}
which is the simulator's $\alpha=\operatorname{atan2}(v_{rel,z},v_{rel,x})$ of
\autoref{sec:relwind} written in ground-frame terms. With wind, $\theta_v$ in
\autoref{eq:alpha_closure} is the air-relative flight-path angle.

\paragraph{The point-mass limit.} Strip the airframe of its rotational degree of freedom: no
aerodynamic normal force, thrust along the velocity vector ($\theta_t=0$ and $\alpha=0$, so
$\hat{x}_v=\hat{x}_r$), a flat non-rotating Earth ($\hat{z}_g=\hat{z}_0$, $\varphi=0$), still air
(so the path speed is the airspeed, $V=v_{rel}$) and drag $D$
antiparallel to $\vec{v}_r$. Then \autoref{eq:eom_trans} reduces to
$m\vec{a}_r = mg\hat{z}_0 + (T-D)\hat{x}_v$, while differentiating
\autoref{eq:vel_axes_in_ground} gives
\begin{equation}\label{eq:accel_in_vel_axes}
    \vec{a}_r = \dot{V}\hat{x}_v - V\dot\theta_v\,\hat{z}_v .
\end{equation}
Projecting on $\hat{z}_v$, where $\hat{z}_0\cdot\hat{z}_v=\cos\theta_v$ and weight is the only
contributor, and on $\hat{x}_v$:
\begin{gather}
    \dot\theta_v = -\frac{g}{V}\cos\theta_v, \label{eq:flightpath_rate}\\
    m\dot{V} = T - D - mg\sin\theta_v. \label{eq:speed_rate}
\end{gather}
\autoref{eq:flightpath_rate} is the \emph{gravity turn}: gravity alone bends the velocity vector, at
a rate $g/V$ modulated by $\cos\theta_v$. It vanishes when the flight is vertical
($\theta_v=\pm90^\circ$), where gravity is purely decelerating and turns nothing, and is largest
when level, $\dot\theta_v=-g/V$. A faster rocket turns less per unit time --- the reason a high
thrust-to-weight launch holds its initial attitude. Holding $V$ fixed, \autoref{eq:flightpath_rate}
integrates to
\begin{equation}\label{eq:gravity_turn_closed}
    \tan\!\left(\frac{\pi}{4}+\frac{\theta_v}{2}\right)
        = \tan\!\left(\frac{\pi}{4}+\frac{\theta_{v,0}}{2}\right) e^{-gt/V},
\end{equation}
so $\theta_v\to-90^\circ$ with time constant $V/g$.

This is a \emph{cruder} model than \autoref{eq:eom_trans}--\autoref{eq:eom_rot}, not a refinement of
them: it carries no moment, no inertia and no angle of attack, and the simulator never integrates it.
Its uses are the two below --- a closed-form estimate to check the integrator against, and the
justification of the short-period approximation that follows.

\subsubsection{Dynamic stability}\label{sec:analysis_dynamic}
The static criterion does not say what the resulting oscillation does. That needs the
\textbf{pitch-damping derivative} $C_{m,q}$, the sensitivity of $C_m$ to the non-dimensional pitch
rate $\hat{\omega}=\omega_r c/(2v_{rel})$:
\begin{equation}
    C_m = C_{m,0} + \alpha\,C_{m,\alpha} + \hat{\omega}\,C_{m,q}.
\end{equation}
It comes from the body's own rotation. A point a distance $s_i$ aft of the CoG moves relative to
the CoG with velocity $\omega_r\hat{y}_r\times\vec{s}_i$ and sees a local angle of attack shifted
by $\Delta\alpha_i\approx+\omega_r s_i/v_{rel}$. Fins well aft see the largest shift, and the
resulting normal-force change against its arm gives a moment proportional to $\omega_r$, present
even at $\alpha=0$.

With the sign of $\vec{A}_m$ used here, a positive $C_m$ pitches the nose \emph{down}. The simulator's
code flags this non-standard sign by renaming the internal quantity to \texttt{C\_m\_code}, keeping the
plain name \texttt{C\_m} reserved at the code's I/O boundary (config files, CSV aero-table headers) for
the standard convention when a user declares \texttt{cm\_convention: standard}. A statically stable,
aft-finned airframe therefore has $C_{m,q}>0$: a positive (nose-up) pitch rate
produces an additional positive (nose-down) moment, opposing the rotation.

Linearising \autoref{eq:moment_aligned} with this term (and $C_{z,0}=C_{m,0}=0$),
\begin{equation}
    I\varepsilon_r = qS(l_g-l_p)C_{z,\alpha}\,\alpha - qSc\,C_{m,\alpha}\,\alpha
        - \frac{qSc^2}{2v_{rel}}C_{m,q}\,\omega_r,
\end{equation}
which separates into a stiffness term and a damping term,
\begin{equation}
    M_\alpha \equiv \frac{\partial\varepsilon_r}{\partial\alpha}
        = \frac{qS}{I}\bigl[(l_g-l_p)C_{z,\alpha}-c\,C_{m,\alpha}\bigr], \qquad
    M_q \equiv \frac{qSc^2}{2v_{rel}I}C_{m,q}.
\end{equation}
$M_\alpha$ is negative for a statically stable airframe and $M_q$ is positive for $C_{m,q}>0$.
Eliminating $\theta_r$ in favour of $\alpha$ through \autoref{eq:alpha_closure}
($\omega_r=\dot\alpha+\dot\theta_v$, $\varepsilon_r=\ddot\alpha+\ddot\theta_v$) turns
$\varepsilon_r=M_\alpha\alpha-M_q\omega_r$ into an exact statement,
\begin{equation}\label{eq:short_period_exact}
    \ddot\alpha + M_q\dot\alpha - M_\alpha\alpha
        = -\left(\ddot\theta_v + M_q\dot\theta_v\right),
\end{equation}
whose right-hand side is the entire content of the short-period approximation: dropping it gives
\begin{equation}
    \frac{\mathrm{d}^2\alpha}{\mathrm{d}t^2} + M_q\frac{\mathrm{d}\alpha}{\mathrm{d}t} - M_\alpha\alpha = 0,
\end{equation}
with natural frequency and damping ratio
\begin{equation}
    \omega_n = \sqrt{-M_\alpha}, \qquad \zeta = \frac{M_q}{2\omega_n} = \frac{M_q}{2\sqrt{-M_\alpha}},
\end{equation}
valid only for $M_\alpha<0$; otherwise the disturbance grows rather than oscillates.

\autoref{eq:flightpath_rate} makes the dropped term quantitative instead of merely ``slow''. With
$|\dot\theta_v|\le g/V$ and $\ddot\theta_v=O\!\left((g/V)^2+g\dot V/V^2\right)$, the dominant
neglected term is $M_q\dot\theta_v=2\zeta\omega_n\dot\theta_v$. Measuring it against the retained
stiffness $M_\alpha\alpha\sim\omega_n^2\hat\alpha$ for an oscillation of amplitude $\hat\alpha$, and
writing
\begin{equation}\label{eq:flightpath_ratio}
    \Lambda \equiv \frac{g/V}{\omega_n} = \frac{g}{V\sqrt{-M_\alpha}},
\end{equation}
the approximation holds when $2\zeta\Lambda\ll\hat\alpha$ and $\Lambda^2\ll\hat\alpha$ --- in short,
$\Lambda\ll1$: the flight-path angle must drift through much less than the oscillation amplitude in
one short period, $\Delta\theta_v\approx2\pi\Lambda\ll\hat\alpha$. Since $\Lambda\propto1/V$, the
approximation is at its weakest at low airspeed --- on and just off the rail, and near the apogee of
a lofted flight --- and strongest through the high-$q$ part of the boost. Near-vertical flight
relaxes it independently of speed, because $\cos\theta_v\to0$ suppresses $\dot\theta_v$ directly.

The distributed-force argument above gives an estimate. A component $i$ with normal-force slope
$C_{z,\alpha,i}$ (referred to $S$) at a distance $s_i$ from the CoG adds a moment proportional to
$\omega_r s_i^2/v_{rel}$, and equating the sum with the $C_{m,q}$ term gives
\begin{equation}
    C_{m,q} \approx \frac{2}{c^2}\sum_i C_{z,\alpha,i}\,s_i^2 ,
\end{equation}
positive in this report's sign convention. OpenRocket uses a different estimate, from the cross-flow
drag of the rotating body and fins \cite{ortech}.
\restodo{a value of $C_{m,q}$ for the project's airframe: the component slopes $C_{z,\alpha,i}$, or
a measured value (forced-oscillation tests or unsteady CFD)}
